\pg{356} b) \emph{Mudança instantanea} ou \emph{successiva}. --- São duas especies da mudança \emph{intrinseca}. --- É \emph{instantanea}, quando se effectua \emph{num instante}, e verifica-se quando o sujeito passa de um para outro estado \emph{contradictorio}, pois as coisas \emph{contradictorias} não admittem meios termos; tal é a mudança de \emph{não-homem} para \emph{homem}. --- É \emph{successiva}, quando se realisa \emph{successivamente}, e dá-se quando o sujeito passa de um para outro estado \emph{contrario}, pois as coisas \emph{contrarias} admittem varios meios termos; tal é a mudança da agua de fria para quente \footnote{A mudança \emph{successiva} chama-se propriamente \emph{movimento}, emquanto a \emph{instantanea} conserva o nome de \emph{mudança}. --- Embora a \emph{mudança} e o \emph{movimento} se empreguem indifferentemente, comtudo differem uma do outro. Com effeito, --- ao passo que o \emph{movimento} exige extremos \emph{positivamente} contrarios, a \emph{mudança} pode effectuar-se entre um extremo negativo e outro positivo, p. ex., entre a \emph{privação} e a \emph{forma}; --- ao passo que o sujeito do \emph{movimento} é um ente dotado de ser substancial, o sujeito da \emph{mudança} pode ser um ente privado do ser substancial, como é a \emph{materia prima}.}.

c) \emph{Mudança substancial} ou \emph{accidental}. --- São duas outras especies da mudança \emph{intrinseca}. --- É \emph{substancial}, quando o sujeito passa de uma para outra \emph{natureza}, ou \emph{essencia}; tal é a mudança do alimento, que se torna carne. --- É \emph{accidental}, quando o sujeito, conservando a sua natureza, passa de uma para outra \emph{propriedade accidental}; tal é a mudança do homem, que de ignorante se torna sabio.

\textbf{§5. O ente, que é acto puro, é immutavel; o ente, que é immutavel, é acto puro.}\footnote{A mudança \emph{substancial} divide-se em \emph{geração} e \emph{corrupção}. --- \emph{A geração} é a passagem do não-ser substancial para o ser substancial, como quando do não-homem se faz o homem. A \emph{corrupção} é a passagem do ser substancial para o não-ser substancial, como quando do homem se faz o cadaver. --- A mudança \emph{substancial} pode dar-se --- ou somente quanto á \emph{forma substancial}, permanecendo a materia, e assim são todas as mudanças naturaes, --- ou quanto á forma e á materia, de modo que \emph{toda} a substancia se mude noutra, e temos a \emph{transubstanciação}. Exemplo d'esta total mudança de uma substancia noutra não se encontra na philosophia, mas unicamente na Theologia revelada, a qual nos ensina que, no S. Sacrificio da Missa, toda a substancia do pão se muda no Corpo adoravel e toda a substancia do vinho se converte no Sangue precioso de N. S. Jesus Cristo. --- A mudança \emph{accidental} chama-se \emph{alteração}, e é a passagem do sujeito para um ser accidental, como quando um homem de ignorante se torna sabio.}
